\chapter{Solution to exercises in chapter 2}
\begin{enumerate}
\item\label{c2p1} Given that the sets $\{\psi_i(\vec{r})\}$ and 
$\{\phi_i(\vec{r})\}$ of $K$ spatial orbitals are orthonormal among themselves 
but not with each other. That is, $\braket{\psi_i|\psi_j} = 
\delta_{ij}, \braket{\phi_i|\phi_j} = \delta_{ij}$ but
\[
\braket{\psi_i|\phi_j} = S_{ij},
\]
which need not be zero. $S$ is called the overlap matrix. Let
\begin{eqnarray*}
\chi_{2i-1}(\vec{x}) &=& \psi_i(\vec{r})\alpha(\omega) \\
\chi_{2i}(\vec{x}) &=& \phi_i(\vec{r})\beta(\omega) \\
\end{eqnarray*}
where $\alpha(+) = 1, \alpha(-) = 0, \beta(+) = 0, \beta(-) = 1$. Consider
\begin{equation}\label{c2pe1}
\braket{\chi_l|\chi_m} = \begin{dcases}
1 & \text{ l = m } \\
\braket{\psi_l|\psi_m} = \delta_{lm} & l \ne m, \text{both odd} \\
\braket{\phi_l|\phi_m} = \delta_{lm} & l \ne m, \text{both even} \\
\braket{\psi_l|\phi_m}\braket{\alpha|\beta} = 0 & l\text{ odd }m\text{ even} \\
\braket{\phi_l|\psi_m}\braket{\beta|\alpha} = 0 & l\text{ even }m\text{ odd}
\end{dcases}
\end{equation}

\item\label{c2p2} Independent electron approximation.
\begin{eqnarray*}
H\Psi^{HP} &=& \sum_{i=1}^N h_i\Psi^{HP} \\
 &=& \sum_{i=1}^N h_i\prod_{j=1}^N\chi_j(\vec{x}_j) \\
 &=& \sum_{i=1}^N \chi_1(\vec{x}_1)\ldots (\varepsilon_i\chi_i(\vec{x}_i))\ldots
 \chi_N(\vec{x_N}) \\
 &=& \left(\sum_{i=1}^N\varepsilon_i\right)\Psi^{HP}
\end{eqnarray*}
Thus $\Psi^{HP}$ is an eigenfunction of $H$ with eigenvalue $\varepsilon_1 + 
\cdots + \varepsilon_N$.

\item\label{c2p3} Normalisation of $2 \times 2$ Slater determinant.
Given that
\[
\Psi(\vec{x}_1, \vec{x}_2) = \frac{1}{\sqrt{2}}
(\chi_1(\vec{x}_1)\chi_2(\vec{x}_2) -\chi_1(\vec{x}_2)\chi_2(\vec{x}_1))
\]
\begin{eqnarray*}
|\Psi|^2 &=& \frac{1}{2}\left[|\chi_1(\vec{x}_1)|^2|\chi_2(\vec{x}_2)|^2 -
\chi_1^\ast(\vec{x}_1)\chi_1(\vec{x}_2)\chi_2^\ast(\vec{x}_2)\chi_2(\vec{x}_1)
 - \right. \\
& & \left.
\chi_1^\ast(\vec{x}_2)\chi_1(\vec{x}_1)\chi_2^\ast(\vec{x}_1)\chi_2(\vec{x}_2)
+ |\chi_1(\vec{x}_2)|^2|\chi_2(\vec{x}_1)|^2\right]
\end{eqnarray*}
We use the integrals:
\begin{eqnarray*}
\int d\vec{x}_1|\chi_1(\vec{x}_1)^2| &=& 1 \\
\int d\vec{x}_2|\chi_1(\vec{x}_2)^2| &=& 1 \\
\int d\vec{x}_1 \chi_1^\ast(\vec{x}_1)\chi_2(\vec{x}_1) &=& 0 \\
\int d\vec{x}_2 \chi_1^\ast(\vec{x}_2)\chi_2(\vec{x}_2) &=& 0 
\end{eqnarray*}
to conclude that 
\[
\int d\vec{x}_1 d\vec{x}_2 |\Psi(\vec{x}_1, \vec{x}_2)|^2 = 1.
\]

\item\label{c2p4} Comparison of eigenvalues of Hatree product and Slater
determinant.
Given that
\begin{eqnarray*}
h(1)\chi_i(\vec{x}_1) &=& \varepsilon_1\chi_i(\vec{x}_1) \\
h(2)\chi_i(\vec{x}_2) &=& \varepsilon_2\chi_i(\vec{x}_2) 
\end{eqnarray*}
and $H = h(1) + h(2)$. Then,
\begin{eqnarray*}
H\Psi^{HP}_{12}(\vec{x}_1, \vec{x}_2) &=& H\chi_i(\vec{x}_1)\chi_j(\vec{x}_2) \\
 &=& h(1)\chi_1(\vec{x}_1)\chi_j(\vec{x}_2) +
     h(2)\chi_1(\vec{x}_1)\chi_j(\vec{x}_2) \\
 &=& \varepsilon_i\chi_1(\vec{x}_1)\chi_j(\vec{x}_2) +
     \varepsilon_j\chi_1(\vec{x}_1)\chi_j(\vec{x}_2) \\
 &=& (\varepsilon_i + \varepsilon_j)\chi_1(\vec{x}_1)\chi_j(\vec{x}_2)
\end{eqnarray*}
Similarly, $H\Psi^{HP}_{21} = (\varepsilon_i + \varepsilon_j)\Psi^{HP}_{21}$.
Now consider
\begin{eqnarray*}
h(1)\Psi(\vec{x}_1, \vec{x}_2) &=& \frac{1}{\sqrt{2}}
(\varepsilon_i\chi_i(\vec{x}_1)\chi_j(\vec{x}_2) - 
\varepsilon_j\chi_j(\vec{x}_1)\chi_i(\vec{x}_2)) \\
h(2)\Psi(\vec{x}_1, \vec{x}_2) &=& \frac{1}{\sqrt{2}}
(\varepsilon_j\chi_i(\vec{x}_1)\chi_j(\vec{x}_2) - 
\varepsilon_i\chi_j(\vec{x}_1)\chi_i(\vec{x}_2)) 
\end{eqnarray*}
so that,
\[
H\Psi(\vec{x}_1, \vec{x}_2) = (\varepsilon_i + \varepsilon_j)
\Psi(\vec{x}_1, \vec{x}_2)
\]
Thus, the Hartee product and the Slater determinant both have the same 
eigenvalues.

\item\label{c2p5} Orthonormality of Slater determinants.
Let $\ket{K} = \ket{\chi_i\chi_j}$ and $\ket{L} = \ket{\chi_k\chi_l}$. Therefore,
\begin{eqnarray*}
\ket{K} &=&\frac{1}{\sqrt{2}}(\chi_i(1)\chi_j(2) - \chi_j(1)\chi_i(2)) \\
\ket{L} &=&\frac{1}{\sqrt{2}}(\chi_k(1)\chi_l(2) - \chi_l(1)\chi_k(2)) 
\end{eqnarray*}
so that
\begin{eqnarray*}
2\braket{K|L} &=& \iint d1d2\left[\chi_i^\ast(1)\chi_k(1)\chi_j^\ast(2)\chi_l(2) - 
 \chi_i^\ast(1)\chi_l(1)\chi_j^\ast(2)\chi_k(2) - \right. \\
 & & \left. \chi_j^\ast(1)\chi_k(1)\chi_i^\ast(2)\chi_l(2) + 
 \chi_j^\ast(1)\chi_l(1)\chi_i^\ast(2)\chi_k(2)\right] \\
 &=& \delta_{ij}\delta_{jl} - \delta_{il}\delta{jk} - \delta_{jk}\delta_{il}
 + \delta_{jl}\delta{ij} \\
\braket{K|L} &=& \delta_{ij}\delta_{jl} - \delta_{il}\delta{jk}
\end{eqnarray*}

\item\label{c2p6} Orthonormality of molecular orbitals.
We will first show that the Guassian orbitals
\[
\phi(\vec{r}) = 
\left(\frac{2\alpha}{\pi}\right)^{3/4}e^{-\alpha|\vec{r}-\vec{R}|^2}
\]
is normalised.
Let
\[
I = \int d\vec{r} |\phi(\vec{r})|^2 = \left(\frac{2\alpha}{\pi}\right)^{3/2}
\iiint dxdydz e^{-2\alpha (\vec{r} - \vec{R})^2}
\]
Let $\vec{r}^\op = \vec{r} - \vec{R}$ so that
\[
I = \left(\frac{2\alpha}{\pi}\right)^{3/2}
\iiint dx^\op dy^\op dz^\op e^{-2\alpha {r^\op}^2}
= \left(\frac{2\alpha}{\pi}\right)^{3/2}
\left(\int_{-\infty}^\infty e^{-2\alpha x^2}dx\right)^3
\]
so that $\braket{\phi|\phi} = 1$.

The remaining part is easy. We have
\begin{eqnarray*}
\psi_1 &=& \frac{\phi_1 + \phi_2}{\sqrt{2(1 + S_{12})}} \\
\psi_2 &=& \frac{\phi_1 - \phi_2}{\sqrt{2(1 - S_{12})}} 
\end{eqnarray*}
where 
\[
S_{12} = \braket{\phi_1|\phi_2}
\]
and
\begin{eqnarray*}
\phi_1&=&\left(\frac{2\alpha}{\pi}\right)^{3/4}e^{-\alpha|\vec{r}-\vec{R}_1|^2} \\
\phi_2&=&\left(\frac{2\alpha}{\pi}\right)^{3/4}e^{-\alpha|\vec{r}-\vec{R}_2|^2},
\end{eqnarray*}
where $\vec{R}_1$ and $\vec{R}_2$ are the positions of the two nuclei.
Now,
\[
\braket{\psi_1|\psi_2} = \frac{\braket{\phi_1|\phi_1} - \braket{\phi_1|\phi_2}
+ \braket{\phi_2|\phi_1} - \braket{\phi_2|\phi_2}}{2\sqrt{1 - S_{12}^2}}
\]
Since the basis functions are real, $\braket{\phi_1|\phi_2} = \braket{\phi_2|
\phi_2}$. They are also normalised so that the numerator vanishes making
$\braket{\psi_1|\psi_2} = 0$.

We next consider,
\[
\braket{\psi_1|\psi_1} = \frac{\braket{\phi_1|\phi_1} + \braket{\phi_1|\phi_2}
+ \braket{\phi_2|\phi_1} + \braket{\phi_2|\phi_2}}{2(1 + S_{12})}
= \frac{1 + S_{12} + S_{12} + 1}{2(1 + S_{12})} = 1
\]
and
\[
\braket{\psi_2|\psi_2} = \frac{\braket{\phi_1|\phi_1} - \braket{\phi_1|\phi_2}
- \braket{\phi_2|\phi_1} + \braket{\phi_2|\phi_2}}{2(1 - S_{12})}
= \frac{1 - S_{12} - S_{12} + 1}{2(1 - S_{12})} = 1
\]
\end{enumerate}
